% !TEX root = ../../../main.tex
% Block card: detection.matrix_profile -- "Matrix profile" (Part I, the blocks chapter).
\cardfield{What it does}
For every window of the signal, it measures how far that window is from its closest
look-alike anywhere else in the span. That list of distances is the
\term{matrix-profile}{matrix profile}.

\begin{plainterms}[The matrix profile]
Cut the trace into every possible stretch of one fixed length, say ten minutes, one starting at
each sample. Take one stretch and slide it along the whole recording, looking for the stretch
whose \emph{shape} is most like it; write down how unlike they still are. Do that for every
stretch. A low number means ``this shape happens again somewhere''; a high number means ``nothing
else looks like this'': a one-off. It is like checking every sentence of a book against every
other sentence and noting how close its nearest echo is. Only the shape counts: each stretch is
shifted and scaled to the same average and spread before the comparison, so a loud and a quiet
copy of the same wiggle count as the same.
\end{plainterms}

\cardfield{When to use it}
When you want to know which parts of a recording repeat and which are unusual, without saying in
advance what you are looking for. It is the foundation of the matrix-profile family: follow it
with \blockref{detection.mp_motifs} to list the shapes that repeat, or with
\blockref{detection.threshold} to mark the stretches with no look-alike.
\blockref{detection.fluss}, \blockref{detection.pan_matrix_profile} and
\blockref{catalogue.snippets} compute their own profiles and do not need this block before them.

\cardfield{What goes in, what comes out}
In goes a span of one channel. Out comes one number per sample: the distance from the window
\emph{starting} at that sample to its nearest look-alike. The last stretch of the span, one window
long less a sample, has no number, because no full window starts there. The profile is also saved
to disk so the matrix-profile pages can show it again without recomputing it.

\cardfield{Settings that matter}
\emph{window\_min} is the window length in minutes, and it is the question you are asking:
``what repeats at this length?''. A pattern much shorter than the window is diluted by everything
around it; one much longer is cut into pieces. \emph{backend} chooses how it is computed; leave it
on \texttt{auto} unless you have a reason.

\cardfield{Watch out for}
The cost grows with the square of the span: twice the span takes four times as long. Long spans
are sent to the cluster or refused (see the entry). A low distance is a matched \emph{shape}, not
a matched size: a stretch of flat drift matches every other stretch of drift almost perfectly, so
remove the baseline first (\blockref{preprocessing.detrend}). Each number belongs to the window
that \emph{starts} at that sample, not to the moment it is drawn at, so an event sits up to a
window to the right of where its low distance is plotted.
